\documentclass[11pt,a4paper]{article}
\usepackage[margin=23mm]{geometry}
\usepackage[T1]{fontenc}
\usepackage[utf8]{inputenc}
\usepackage{lmodern,microtype,amsmath,amssymb,booktabs,tabularx,array}
\usepackage{graphicx,xcolor,listings,enumitem}
\usepackage[section]{placeins}
\usepackage[colorlinks=true,linkcolor=blue!50!black,urlcolor=blue!50!black]{hyperref}
\usepackage{fancyhdr}
\pagestyle{fancy}
\fancyhf{}
\lhead{Statistical learning: worked notebook explanations}
\rfoot{\thepage}
\setlength{\headheight}{14pt}
\setlength{\parindent}{0pt}
\setlength{\parskip}{6pt}
\setlength{\emergencystretch}{2em}
\setlist{itemsep=3pt,topsep=4pt}
\lstset{language=Python,basicstyle=\small\ttfamily,breaklines=true,
 columns=fullflexible,keepspaces=true,showstringspaces=false,
 frame=single,rulecolor=\color{gray!35},backgroundcolor=\color{gray!4},
 commentstyle=\color{green!35!black},keywordstyle=\color{blue!65!black},
 aboveskip=8pt,belowskip=8pt}
\newcommand{\code}[1]{\texttt{\detokenize{#1}}}
\newcommand{\RSS}{\operatorname{RSS}}
\newcommand{\RMSE}{\operatorname{RMSE}}
\newcommand{\MAE}{\operatorname{MAE}}
\newcommand{\E}{\mathbb{E}}
\newcommand{\ind}[1]{\mathbf{1}\{#1\}}
\newcommand{\fig}[2]{\begin{figure}[htbp]\centering
\includegraphics[width=\linewidth,height=0.39\textheight,keepaspectratio]{figures/#1}
\caption{#2}\end{figure}}
\newcommand{\notebooktitle}[3]{\begin{center}{\LARGE\bfseries #1\par}\vspace{5pt}
{\large Problem statement and worked resolution\par}\vspace{5pt}
Companion to \code{#2}\par\end{center}
\textbf{How to use this explanation.} #3 Section numbers follow the notebook's numbered sections where present. Numerical values are rounded and refer to the recorded run, not universal performance claims.}
\newcommand{\methods}[1]{\subsection*{Main Python calls}#1}

\begin{document}
\notebooktitle{A: Request latency, least squares, and ridge}{A_request_latency.ipynb}{Read this alongside the executable cells. The resolution explains both the fitting algorithms and why stable predictions do not always require stable individual coefficients.}

\section*{Problem and notation}
Predict request latency $y$ in milliseconds from payload $s$ in MB and queue length $q$.
The simulator generates 600 independent requests with
\[
s\sim U(0.5,10),\quad q\sim\operatorname{Poisson}(3),\quad
 y=40+8s+3q+\epsilon,\quad \epsilon\sim N(0,12^2).
\]
All draws above are independent. A second payload measurement is $s'=s+u$, with independent $u\sim N(0,0.02^2)$. The prime labels a noisy copy, not a derivative.
The questions are: can we recover the linear relationship, verify different optimizers, and stabilize coefficients when two inputs carry almost identical information?

An intercept $b$ and slope vector $w$ define $\hat y=b+x^Tw$.
The vector $\beta=(b,w^T)^T$ includes the intercept. In matrix formulas $X$ includes a leading ones column, so $\hat y=X\beta$; scikit-learn receives feature columns only and handles its own intercept.
Ordinary least squares (OLS) minimizes total squared residuals. A residual is $y_i-\hat y_i$.
Cross-validation (CV) scores a model on training rows excluded from that particular fit.

\section{Generate data and reserve partitions}
Use simulation seed 17 and split seed 23. Reserve 120 test rows, then 120 validation rows from the remaining 480; 360 rows remain for training. Validation provides an audit, while five-fold CV within training selects ridge strength. Test rows are used only at the end.

\methods{\code{np.random.default_rng} supplies reproducible random draws; \code{train_test_split} separates rows; \code{DataFrame} preserves feature names and row identifiers. Index-disjointness assertions check partition separation.}
The duplicated feature is generated without learning anything from targets. Generating it before splitting is safe here; estimating a mean, scale, or model from all rows would not be.
The training scatterplot shows the payload relationship, with spread caused by queue variation and outcome noise. It does not reveal an exact deterministic line.

\section{Fit a constant and a payload-only line}
For a constant prediction, minimizing $\sum_i(y_i-b)^2$ gives $b=\bar y$.
For the line $b+ws$, differentiate with respect to $b$ and $w$ and solve:
\[
w=\frac{\sum_i(s_i-\bar s)(y_i-\bar y)}{\sum_i(s_i-\bar s)^2},
\qquad b=\bar y-w\bar s.
\]
The denominator is nonzero because payload varies. The implementation computes this formula and compares it with \code{LinearRegression().fit(...)} using \code{assert_allclose}.

The recorded baseline is $88.998$ ms; the fitted line is
\[\hat y=45.849+8.524s.\]
The slope means an additional predicted $8.524$ ms per MB. It need not equal the generating slope 8 in a finite sample. Since queue is omitted, the intercept also absorbs its average contribution; it should not simply be compared with 40 as though the models had the same inputs. In the population, independence gives $\E[y\mid s]=49+8s$.

\section{Fit both features using matrices}
With columns $(1,s,q)$, solve
\[\min_\beta\|X\beta-y\|^2.\]
Setting the gradient to zero yields the normal equations $X^TX\beta=X^Ty$.
These equations explain the solution, but explicitly forming $(X^TX)^{-1}$ is an unnecessary and potentially unstable computational step. The notebook uses \code{np.linalg.lstsq} and verifies agreement with \code{LinearRegression}.
\begin{center}
\begin{tabular}{lrr}\toprule
Parameter & Generating value & Fitted value\\\midrule
Intercept (ms) &40&36.2358\\
Payload (ms/MB)&8&8.4116\\
Queue (ms/request)&3&3.6336\\\bottomrule
\end{tabular}
\end{center}
Each slope now describes a change with the other input held fixed. The numerical solvers agree; discrepancies from generating coefficients are statistical estimation error, not evidence that a solver failed.

\section{Solve the same problem by gradient descent}
For $n$ fitting rows, define
\[
J(\beta)=\frac{1}{2n}\|X\beta-y\|^2,\qquad
\nabla J(\beta)=\frac1nX^T(X\beta-y).
\]
Starting at zero, update $\beta\leftarrow\beta-\eta\nabla J$. The Hessian is $X^TX/n$, so a safe step is $\eta=1/L$ with $L=\|X\|_2^2/n$; the matrix norm is the largest singular value. Stop when the gradient norm is below $10^{-8}$, with a 20,000-update limit.

\methods{\code{least_squares_gd} is the notebook's implementation; \code{np.linalg.norm(..., ord=2)} computes the spectral norm. \code{StandardScaler.fit} learns training means $\mu_j$ and standard deviations $a_j$; \code{transform} applies $z_j=(x_j-\mu_j)/a_j$.}
Scaling improves numerical conditioning without changing the set of linear prediction functions. If the standardized coefficients are $(\tilde b,\tilde w)$, recover original-unit coefficients using
\[w_j=\tilde w_j/a_j,\qquad b=\tilde b-\sum_jw_j\mu_j.\]
Both runs agree with least squares within roughly $10^{-7}$ ms in training predictions. Raw features require 5,919 updates, whereas standardized features require 9. These are iteration counts for the stated stopping criterion, not a universal speed ratio.
\fig{A_request_latency_3.png}{The two optimizations reach the same objective. Standardization changes the geometry and greatly reduces the update count in this example.}

\section{Choose ridge strength on the expanded features}
With $(s,q,s')$, two columns are nearly equal. Changes to the payload coefficients can offset each other, producing little prediction change. This makes individual OLS slopes poorly determined.
Ridge instead minimizes
\[
\frac{1}{2n}\|X\beta-y\|^2+\frac\lambda2\|w\|^2,
\]
where $w$ contains the three standardized slopes and excludes the intercept. The penalty discourages large compensating coefficients. Scaling matters for ridge because the penalty depends on coefficient units.

\methods{\code{KFold} supplies five fitting/scoring partitions; \code{Ridge(alpha=...)} fits penalized least squares. Its objective is RSS plus \code{alpha} times the squared slope norm, so multiplying our objective by $2n$ gives \code{alpha = n_fit * lambda}.}
Search $\lambda\in\{10^{-4},10^{-3},10^{-2},10^{-1},1,10\}$. Each fold fits on 288 rows and scores on 72. Fit a fresh scaler on the 288 rows, transform both subsets with that scaler, fit ridge, and score held-out RMSE. Reuse the same folds for all penalties. Average fold RMSE values; an exact tie uses grid order. Fold standard deviation describes score variability, not a confidence interval.

The selected $\lambda$ is $10^{-4}$, with CV RMSE about $12.183$ ms. Very strong shrinkage, $\lambda=10$, gives about $22.886$ ms. Refit the scaler and model on all 360 training rows; the final library penalty is $360\times10^{-4}=0.036$. Also fit expanded OLS as a prespecified comparison. Do not refit on validation after inspecting it.

\section{Inspect shrinkage and bootstrap stability}
The ridge path plots absolute standardized slopes as the penalty varies. Absolute values hide signs, so this plot illustrates magnitudes, not the direction of every association. Under correlated inputs, individual coefficient paths need not decrease monotonically.

A bootstrap sample draws 360 training rows with replacement. Across 100 resamples, refit both scaler and model, convert slopes back to original units, and predict at $(s,q,s')=(5,3,5)$. Keep the chosen $\lambda$ fixed: this assesses fitting variability at that setting, not the entire selection procedure.
\begin{center}
\begin{tabular}{lrr}\toprule
Bootstrap standard deviation & OLS & Ridge\\\midrule
Payload slope (ms/MB)&33.914&7.668\\
Copy slope (ms/MB)&33.930&7.683\\
Sum of payload slopes (ms/MB)&0.235&0.234\\
Fixed-request prediction (ms)&0.641&0.630\\\bottomrule
\end{tabular}
\end{center}
The individual coefficients fluctuate substantially, but their sum is stable. Indeed,
\[w_s s+w_{s'}s'=(w_s+w_{s'})s+w_{s'}u.\]
When $u$ is small, predictions largely depend on the sum. Ridge reduces individual coefficient variability without dramatically changing prediction variability at this typical input. This does not justify prediction for inputs with $s'$ far from $s$.
\fig{A_request_latency_5.png}{Large opposite changes in the two OLS payload slopes can cancel. The fixed-input prediction distributions are much tighter than the separate coefficient distributions suggest.}

\section{Residuals and validation audit}
Training residuals from the two-feature model should have no systematic trend if the mean relationship is well captured. A curved pattern can indicate a missing relationship; changing spread can indicate varying noise. Training residuals are optimistic for evaluating new-row prediction.
Queue and payload are independent in this generator, so omitting queue need not produce an obvious residual trend against payload. A lack of visible trend is not proof that all useful inputs have been included.
\code{regression_report} predicts with each frozen model and computes mean absolute error (MAE) and RMSE:
\[\MAE=\frac1m\sum_i|y_i-\hat y_i|,\qquad
\RMSE=\sqrt{\frac1m\sum_i(y_i-\hat y_i)^2}.\]
Both are in milliseconds. The validation audit leaves all choices unchanged; it gives approximately $12.505$ ms RMSE for the two-feature model and $12.519$ for expanded ridge.

\section{Final test and answers to the interpretation questions}
\begin{center}
\begin{tabular}{lrr}\toprule
Frozen model & Test MAE (ms)&Test RMSE (ms)\\\midrule
Training mean&21.702&25.810\\
Payload only&10.427&12.466\\
Payload plus queue&10.299&12.737\\
Expanded OLS&10.321&12.738\\
Expanded ridge&10.300&12.729\\\bottomrule
\end{tabular}
\end{center}
The simple payload-only model happens to have the lowest test RMSE, while the two-feature model has lower MAE. Different losses weight errors differently, and a 120-row test sample need not reproduce population rankings. Do not use this table to repeat model selection.

The key resolutions are: all three numerical approaches agree; standardization improves optimization; collinearity can destabilize coefficients more than predictions; and ridge stabilizes the duplicated slopes without guaranteeing a large improvement in held-out error. The independent noise has standard deviation 12 ms, so even the exact conditional mean has population RMSE 12 ms. A finite test score may lie above or below it.

\section*{Further explanation: why these results occur}
\subsection*{The geometry of a nearly duplicated feature}
Imagine first that $s'=s$ exactly. Then
$w_s s+w_{s'}s'=(w_s+w_{s'})s$: the data identify the sum, but cannot distinguish
$(w_s,w_{s'})=(8,0)$ from $(108,-100)$. This is a lack of identifiability, not a failure of Python.
With a small difference $u=s'-s$, these alternatives become distinguishable only through the tiny term $w_{s'}u$.
Noise in the targets can then produce large changes in the two separate fitted slopes.

Ridge makes an additional preference explicit. For a fixed sum $a=w_s+w_{s'}$, the identity
\[
w_s^2+w_{s'}^2=\frac{a^2}{2}+\frac{(w_s-w_{s'})^2}{2}
\]
shows why its penalty prefers balanced coefficients when the two columns are equal and identically scaled.
With a near-copy and a weak selected penalty they need not be equal or even both positive.
Thus the negative mean payload slope in the bootstrap table does not by itself imply that larger payloads reduce latency.

\subsection*{Why the same redundancy also affects optimization}
For centered, standardized feature matrix $Z$, the ridge normal equations are
\[(Z^TZ+n\lambda I)w=Z^T(y-\bar y\mathbf 1).\]
If $Z^TZ$ has a very small eigenvalue $d$, unpenalized fitting divides the corresponding target projection by $d$.
Ridge replaces that divisor by $d+n\lambda$, preventing enormous amplification in that direction.
This explains both numerical and statistical stabilization. It does not remove the underlying ambiguity about separate real-world effects.

For gradient descent on a quadratic, coefficient error in an eigenvector direction is multiplied by
$1-\eta d/n$ per update. When eigenvalues have very different sizes, a step small enough for the largest
one makes slow progress in the smallest direction. Standardization often improves this imbalance,
but it does not generally remove collinearity. Scaling and ridge solve related but different problems.

\subsection*{Prediction uncertainty is not just coefficient uncertainty}
The bootstrap fixed-request distribution measures variability in the estimated conditional mean.
A future observed latency also contains new noise. Conditional on the training data, if the true mean is $f(x)$,
\[
\E[(Y_{\rm new}-\hat f(x))^2\mid x,\text{training}]
=12^2+(f(x)-\hat f(x))^2.
\]
The cross-term vanishes because future noise is independent and mean zero.
Consequently a bootstrap prediction standard deviation near 0.63 ms does not imply that future latencies lie within
roughly 0.63 ms of the prediction. The irreducible noise is much larger.

\subsection*{Documentation details that change how code should be read}
The \href{https://scikit-learn.org/1.8/modules/generated/sklearn.linear_model.LinearRegression.html}{LinearRegression reference}
separates constructor parameters from fitted attributes and methods. In particular, \code{fit_intercept=True}
means the library supplies the intercept, and \code{score} returns $R^2$, not RMSE.
The \href{https://scikit-learn.org/1.8/modules/generated/sklearn.linear_model.Ridge.html}{Ridge reference}
states its summed-loss objective; that is why our average-loss $\lambda$ requires the factor $n$.
The trailing underscore in \code{coef_} means an attribute learned during fitting.

\textbf{Practical checks.} Keep original feature units in your reported coefficients; preserve column order;
reuse the training scaler for predictions; and inspect predictions as well as coefficients when collinearity is present.
A solver tolerance tests optimization accuracy, not correctness of the statistical model.
If the selected penalty lies at the edge of a grid, a wider search may be a sensible \emph{new development experiment},
but the existing final test must not be repeatedly reused as an untouched assessment.

\end{document}
